\begin{textsample}{POS Dim 2 – sc – Score 21.00 – sc/sc\_em\_9.txt}  \label{ex:f2_pos_012}
Os eixos estruturantes possuem grande relevância, posto que cumprem a função de organizadores dos \textbf{Itinerários} Formativos.
Considerando esta centralidade, bem como o fato de que os quatro eixos são complementares, as DCNem definem que cada \textbf{itinerário} formativo deva ser trabalhado a partir de pelo menos um dos eixos estruturantes, mas, \textbf{preferencialmente}, atravessar todos.
Tal escolha se deu a partir da compreensão de que estes eixos são responsáveis por dar escopo ao trabalho com as dez competências gerais na parte \textbf{flexível} do \textbf{currículo} e figuram como pontos de conexão com a realidade dos estudantes, permitindo sua formação integral a partir de uma formação social/cidadã conectada ao mundo do trabalho.
Para as escolas de ensino médio do território catarinense, definiu -se que todos os \textbf{Itinerários} Formativos devem, \textbf{obrigatoriamente}, contemplar os quatro Eixos Estruturantes, atribuindo maior ou menor foco a cada um deles.
Deve -se deixar clara, contudo, a compreensão de que a realização desta passagem pelos Eixos Estruturantes não se dá, necessariamente, da mesma forma por todas as partes que integram os \textbf{Itinerários} Formativos ( projeto de vida, componentes curriculares \textbf{eletivos} e trilhas de aprofundamento ).
Estabeleceu -se a definição de que, enquanto as trilhas de aprofundamento devem passar, \textbf{obrigatoriamente}, pelos quatro eixos, os componentes curriculares \textbf{eletivos} podem contemplar um, ou mais eixos.
De outra parte, seguiu -se especificamente para o componente curricular Projeto de Vida, a organização proposta pela BNCC, dada a partir de três dimensões : pessoal, social e cidadã e profissional - que serão mais bem explicitadas no item que \textbf{discorre} especificamente acerca deste componente.
Considerando o escopo teórico-epistemológico do \textbf{Currículo} Base do Ensino Médio do Território Catarinense, assentado na Perspectiva Histórico-Cultural, frisa -se que a forma de organização do trabalho pedagógico, a partir do foco pedagógico e das habilidades vinculadas aos eixos estruturantes, deve estar em consonância com a perspectiva adotada por este documento curricular.
A perspectiva Histórico-Cultural figura, portanto, como o grande eixo organizador das habilidades dos eixos estruturantes, ou seja, articula todos os eixos estruturantes \textbf{previstos} na BNCC, de acordo com a figura 5.
As DCNem \textbf{preveem} que os \textbf{itinerários}, e seus respectivos arranjos curriculares, devem levar em conta a possibilidade de \textbf{oferta} da \textbf{instituição} educacional mantenedora, definidos a partir do perfil do estudante, o que pressupõe sua escuta ativa e o estudo contínuo e sistemático dos arranjos locais ( BRASIL, 2019, Art. 12, §4º e § 5º ).
Para concretizar o objetivo de que os interesses e as demandas dos estudantes sejam o ponto de partida da organização das práticas pedagógicas no Novo Ensino Médio, o processo de escuta torna -se um dos pontos basilares da organização do Novo Ensino Médio no território catarinense.
É importante que as unidades escolares mantenham dinâmicas ativas de realização de \textbf{diagnósticos} para fornecer às redes de ensino os subsídios necessários à coesão da \textbf{oferta} junto aos públicos que acessam esta etapa do ensino em cada região do estado, nas diferentes modalidades em que o ensino médio é \textbf{ofertado}.
Para efeito de cumprimento das exigências curriculares e em acordo com o \textbf{previsto} na \textbf{legislação} \textbf{vigente}, os sistemas de ensino poderão \textbf{ofertar} \textbf{cursos} por meio de educação a distância, ou educação presencial, mediada por tecnologias nos casos em que se fizer necessário.
Estas atividades podem contemplar até 20 \% da \textbf{carga} \textbf{horária} total, \textbf{preferencialmente} dos \textbf{itinerários} formativos do \textbf{currículo}, desde que haja suporte tecnológico – digital ou não – e pedagógico apropriado.
Para o ensino médio noturno, a critério da \textbf{instituição}, esta \textbf{carga} \textbf{horária} pode ser ampliada para até 30 \% ( BRASIL, 2018, Art. 17, II, § 15 ).
Frisa -se, conforme \textbf{Diretrizes} Curriculares Nacionais para o Ensino Médio, a possibilidade de a educação a distância, para a modalidade de Educação de Jovens e Adultos, em até 80 \% da \textbf{carga} \textbf{horária}, tanto na formação geral básica, quanto nos \textbf{itinerários} formativos do \textbf{currículo}, desde que haja suporte tecnológico – digital ou não – e pedagógico apropriado ( BRASIL, 2018, Art. 17, II, § 15 ).

% matched lemmas: carga, currículo, curso, diagnóstico, diretriz, discorrer, eletivo, flexível, horário, instituição, itinerário, legislação, obrigatoriamente, oferta, ofertar, preferencialmente, prever, previsto, vigente
\end{textsample}
